\documentclass[12pt,a4paper]{article}

\usepackage[utf8]{inputenc}
\usepackage[T1]{fontenc}
\usepackage[french]{babel}
\usepackage{geometry}
\usepackage{booktabs}
\usepackage{xcolor}
\usepackage{listings}
\usepackage{hyperref}
\usepackage{titlesec}
\usepackage{enumitem}
\usepackage{fancyhdr}
\usepackage{array}
\usepackage{tabularx}
\usepackage{mdframed}

\geometry{margin=2.5cm}

% Couleurs
\definecolor{primary}{RGB}{0, 102, 204}
\definecolor{success}{RGB}{40, 167, 69}
\definecolor{danger}{RGB}{220, 53, 69}
\definecolor{warning}{RGB}{255, 193, 7}
\definecolor{codebg}{RGB}{245, 245, 245}
\definecolor{codeframe}{RGB}{200, 200, 200}

% Style des listings
\lstset{
    backgroundcolor=\color{codebg},
    basicstyle=\ttfamily\small,
    breaklines=true,
    frame=single,
    rulecolor=\color{codeframe},
    showstringspaces=false,
    keywordstyle=\color{primary}\bfseries,
    commentstyle=\color{gray}\itshape,
    stringstyle=\color{success},
    language=PHP,
    literate=
        {é}{{\'{e}}}1 {è}{{\`{e}}}1 {ê}{{\^{e}}}1
        {à}{{\`{a}}}1 {â}{{\^{a}}}1 {î}{{\^{i}}}1
        {ô}{{\^{o}}}1 {ù}{{\`{u}}}1 {û}{{\^{u}}}1
        {ç}{{\c{c}}}1 {œ}{{\oe}}1 {æ}{{\ae}}1
}

% En-têtes et pieds de page
\pagestyle{fancy}
\fancyhf{}
\rhead{\textcolor{primary}{\textbf{EILCO — Vie Associative}}}
\lhead{Rapport de corrections}
\rfoot{\thepage}
\lfoot{\today}

% Style des titres
\titleformat{\section}{\large\bfseries\color{primary}}{}{0em}{}[\titlerule]
\titleformat{\subsection}{\normalsize\bfseries\color{primary!80!black}}{}{0em}{}

% Boîte de code
\newenvironment{codebox}[1]{%
    \begin{mdframed}[backgroundcolor=codebg, linecolor=codeframe, linewidth=1pt]%
    \textbf{\texttt{#1}}\\[4pt]%
}{%
    \end{mdframed}%
}

\begin{document}

% ============================================================
% PAGE DE TITRE
% ============================================================
\begin{titlepage}
    \centering
    \vspace*{2cm}

    {\Huge\bfseries\color{primary} Rapport de Corrections}\\[0.5cm]
    {\Large\color{gray} Projet Vie Associative — EILCO}\\[2cm]

    \begin{mdframed}[backgroundcolor=primary!10, linecolor=primary, linewidth=2pt]
        \centering
        \vspace{0.5cm}
        {\large\bfseries Résumé exécutif}\\[0.5cm]
        \normalsize
        Ce rapport documente l'ensemble des corrections de bogues et des
        améliorations d'expérience utilisateur apportées à l'application
        \textit{Vie Associative EILCO} lors de la session du
        \textbf{28 mars 2026}.\\[0.5cm]
        \textbf{5 corrections} ont été apportées sur \textbf{5 fichiers distincts}.
        \vspace{0.3cm}
    \end{mdframed}

    \vspace{2cm}

    \begin{tabular}{ll}
        \textbf{Projet}     & Vie Associative EILCO \\
        \textbf{Branche}    & Youssef \\
        \textbf{Date}       & 28 mars 2026 \\
        \textbf{Technologie} & PHP 8+ (MVC custom), MySQL, JavaScript \\
        \textbf{Serveur}    & XAMPP — \texttt{localhost:8080} \\
    \end{tabular}

    \vfill
    {\small\color{gray} Généré automatiquement par Claude Code — Anthropic}
\end{titlepage}

% ============================================================
% TABLE DES MATIÈRES
% ============================================================
\tableofcontents
\newpage

% ============================================================
\section{Contexte et diagnostic initial}
% ============================================================

L'application ne répondait pas dans le navigateur (page bloquée sur
\textit{Loading\ldots} à l'adresse \texttt{http://localhost:8080/}).
Une analyse en plusieurs étapes a permis d'identifier la chaîne de causes :

\begin{enumerate}[leftmargin=2cm]
    \item \textbf{MySQL non démarré} — le serveur PHP acceptait les connexions
          TCP mais toute requête se bloquait indéfiniment en attendant la base
          de données.
    \item \textbf{Erreur 500 après démarrage de MySQL} — la méthode
          \texttt{Security::getCsrfToken()} appelée dans une vue n'existe pas
          dans la classe \texttt{Security}.
    \item \textbf{Avertissements PHP} — variable \texttt{\$months\_fr}
          non définie dans le contrôleur utilisateur.
    \item \textbf{Problèmes UX mineurs} — titre de page manquant sur le
          tableau de bord admin, lien de navigation menant à une page
          restreinte pour les utilisateurs de permission 1, et URLs de
          connexion/inscription incorrectes.
\end{enumerate}

% ============================================================
\section{Corrections apportées}
% ============================================================

\subsection{Correction 1 — Méthode CSRF inexistante (Bogue critique)}

\begin{mdframed}[backgroundcolor=danger!10, linecolor=danger, linewidth=1.5pt]
\textbf{Sévérité :} Critique — empêche le chargement de toutes les pages\\
\textbf{Fichier :} \texttt{views/includes/calendrier-general.php}, ligne 22
\end{mdframed}

\vspace{0.5cm}
\textbf{Cause :} Le composant calendrier appelait la méthode
\texttt{getCsrfToken()} qui n'existe pas dans la classe \texttt{Security}.
La méthode correcte est \texttt{generateCsrfToken()}.

\textbf{Avant :}
\begin{lstlisting}
data-csrf="<?= htmlspecialchars(Security::getCsrfToken(),
    ENT_QUOTES, 'UTF-8') ?>"
\end{lstlisting}

\textbf{Après :}
\begin{lstlisting}
data-csrf="<?= htmlspecialchars(Security::generateCsrfToken(),
    ENT_QUOTES, 'UTF-8') ?>"
\end{lstlisting}

\textbf{Impact :} Sans cette correction, toute requête vers la page d'accueil
(qui inclut le calendrier) provoquait une erreur fatale PHP, rendant
l'application totalement inaccessible.

% ---
\subsection{Correction 2 — Variable \texttt{\$months\_fr} non définie}

\begin{mdframed}[backgroundcolor=warning!20, linecolor=warning, linewidth=1.5pt]
\textbf{Sévérité :} Haute — avertissement PHP affiché, données manquantes\\
\textbf{Fichier :} \texttt{controllers/UserController.php}, méthode \texttt{dashboard()}
\end{mdframed}

\vspace{0.5cm}
\textbf{Cause :} La vue \texttt{views/user/dashboard.php} utilisait le tableau
\texttt{\$months\_fr} (noms des mois en français) pour afficher les événements
recommandés, mais ce tableau n'était jamais passé par le contrôleur.

\textbf{Avant :}
\begin{lstlisting}
return [
    'user'               => $user,
    'stats'              => $stats,
    'my_clubs'           => $my_clubs,
    'upcoming_events'    => array_slice($upcoming_events, 0, 5),
    'past_events'        => array_slice($past_events, 0, 3),
    'recommended_events' => $recommended_events
];
\end{lstlisting}

\textbf{Après :}
\begin{lstlisting}
$months_fr = ['Jan','Fev','Mar','Avr','Mai','Juin',
              'Juil','Aout','Sep','Oct','Nov','Dec'];

return [
    'user'               => $user,
    'stats'              => $stats,
    'my_clubs'           => $my_clubs,
    'upcoming_events'    => array_slice($upcoming_events, 0, 5),
    'past_events'        => array_slice($past_events, 0, 3),
    'recommended_events' => $recommended_events,
    'months_fr'          => $months_fr
];
\end{lstlisting}

\textbf{Impact :} Sans cette correction, les noms de mois n'apparaissaient pas
dans la section « Recommandés pour vous » du tableau de bord utilisateur, et
PHP générait un avertissement visible dans les logs.

% ---
\subsection{Correction 3 — Titre de page manquant (Admin)}

\begin{mdframed}[backgroundcolor=warning!20, linecolor=warning, linewidth=1.5pt]
\textbf{Sévérité :} Moyenne — UX dégradée\\
\textbf{Fichier :} \texttt{views/admin/dashboard.php}
\end{mdframed}

\vspace{0.5cm}
\textbf{Cause :} La variable \texttt{\$pageTitle} n'était pas définie dans la
vue du tableau de bord administrateur, ce qui laissait l'onglet du navigateur
vide ou avec une valeur par défaut non pertinente.

\textbf{Avant :}
\begin{lstlisting}
$pageCss = ['shared', 'buttons', 'tables', 'admin', 'dashboard'];
\end{lstlisting}

\textbf{Après :}
\begin{lstlisting}
$pageTitle = 'Tableau de bord - Administration EILCO';
$pageCss = ['shared', 'buttons', 'tables', 'admin', 'dashboard'];
\end{lstlisting}

% ---
\subsection{Correction 4 — Lien « Découvrir les clubs » inaccessible}

\begin{mdframed}[backgroundcolor=warning!20, linecolor=warning, linewidth=1.5pt]
\textbf{Sévérité :} Haute — erreur 403 pour les étudiants ordinaires\\
\textbf{Fichier :} \texttt{views/user/dashboard.php}, ligne 55
\end{mdframed}

\vspace{0.5cm}
\textbf{Cause :} Le bouton « Découvrir les clubs » pointait vers
\texttt{?page=club-list}, une route réservée aux utilisateurs de
permission $\geq$ 2 (tuteurs, BDE, admins). Les étudiants classiques
(permission 1) obtenaient une erreur 403 en cliquant dessus.

\textbf{Avant :}
\begin{lstlisting}
<a href="?page=club-list" class="hero-btn">
    <i class="fas fa-users"></i>
    Decouvrir les clubs
</a>
\end{lstlisting}

\textbf{Après :}
\begin{lstlisting}
<?php if (($_SESSION['permission'] ?? 0) >= 2): ?>
<a href="?page=club-list" class="hero-btn">
    <i class="fas fa-users"></i>
    Decouvrir les clubs
</a>
<?php else: ?>
<a href="?page=home" class="hero-btn">
    <i class="fas fa-users"></i>
    Decouvrir les clubs
</a>
<?php endif; ?>
\end{lstlisting}

\textbf{Impact :} Les étudiants (permission 1) sont redirigés vers la page
d'accueil qui présente déjà les clubs par campus, tandis que les utilisateurs
privilégiés accèdent à la liste complète de gestion.

% ---
\subsection{Correction 5 — URLs de connexion/inscription incorrectes}

\begin{mdframed}[backgroundcolor=warning!20, linecolor=warning, linewidth=1.5pt]
\textbf{Sévérité :} Moyenne — liens cassés sur certaines configurations\\
\textbf{Fichier :} \texttt{views/includes/header.php}, lignes 179--184
\end{mdframed}

\vspace{0.5cm}
\textbf{Cause :} Les boutons « Connexion » et « Inscription » dans l'en-tête
utilisaient des URLs avec un slash initial (\texttt{/?page=login}), ce qui
peut causer des erreurs de routage selon la configuration du serveur.
L'application utilise des URLs relatives sans slash initial partout ailleurs.

\textbf{Avant :}
\begin{lstlisting}
<a href="/?page=login" class="header-btn header-btn-outline">
    Connexion
</a>
<a href="/?page=register" class="header-btn header-btn-primary">
    Inscription
</a>
\end{lstlisting}

\textbf{Après :}
\begin{lstlisting}
<a href="?page=login" class="header-btn header-btn-outline">
    Connexion
</a>
<a href="?page=register" class="header-btn header-btn-primary">
    Inscription
</a>
\end{lstlisting}

% ============================================================
\section{Tableau récapitulatif}
% ============================================================

\begin{center}
\renewcommand{\arraystretch}{1.4}
\begin{tabularx}{\textwidth}{
    >{\bfseries}p{0.5cm}
    >{\ttfamily\small}p{4.5cm}
    X
    p{2.2cm}
}
\toprule
\textbf{N°} & \textbf{Fichier} & \textbf{Description} & \textbf{Sévérité} \\
\midrule
1 & views/includes/\\calendrier-general.php
  & Méthode CSRF inexistante (\texttt{getCsrfToken} → \texttt{generateCsrfToken})
  & \textcolor{danger}{\textbf{Critique}} \\
\midrule
2 & controllers/\\UserController.php
  & Variable \texttt{\$months\_fr} manquante dans le retour du contrôleur
  & \textcolor{warning!80!black}{\textbf{Haute}} \\
\midrule
3 & views/admin/\\dashboard.php
  & \texttt{\$pageTitle} absent — onglet navigateur vide
  & \textcolor{primary}{\textbf{Moyenne}} \\
\midrule
4 & views/user/\\dashboard.php
  & Lien « clubs » → erreur 403 pour les étudiants (permission 1)
  & \textcolor{warning!80!black}{\textbf{Haute}} \\
\midrule
5 & views/includes/\\header.php
  & URLs \texttt{/?page=} incorrectes → \texttt{?page=}
  & \textcolor{primary}{\textbf{Moyenne}} \\
\bottomrule
\end{tabularx}
\end{center}

% ============================================================
\section{Note sur l'environnement de développement}
% ============================================================

L'application repose sur \textbf{XAMPP} pour le serveur web (Apache) et la
base de données (MySQL). Lors de la session, MySQL n'était pas démarré, ce
qui provoquait un blocage indéfini de toutes les requêtes PHP.

\begin{mdframed}[backgroundcolor=primary!8, linecolor=primary, linewidth=1pt]
\textbf{Procédure à suivre avant chaque session de développement :}
\begin{enumerate}
    \item Ouvrir le panneau de contrôle XAMPP.
    \item Démarrer les services \textbf{Apache} et \textbf{MySQL}.
    \item Vérifier que \texttt{http://localhost:8080/} répond correctement.
\end{enumerate}
Ou via la ligne de commande : \texttt{C:/xampp/mysql\_start.bat}
\end{mdframed}

% ============================================================
\section{Conclusion}
% ============================================================

Les cinq corrections apportées ont permis de restaurer le bon fonctionnement
de l'application et d'améliorer l'expérience utilisateur pour les étudiants
de faible niveau de permission. La correction la plus critique était
l'appel à \texttt{Security::getCsrfToken()} qui rendait l'intégralité
du site inaccessible.

Aucune régression n'a été introduite : les pages principales
(\texttt{home}, \texttt{login}, \texttt{event-list}) ont été testées
et répondent correctement après les corrections.

\vspace{1cm}
\hrule
\vspace{0.3cm}
{\small\color{gray}
Rapport généré automatiquement — Claude Code (Anthropic) —
Session du 28 mars 2026
}

\end{document}
